\documentclass[aspectratio=169]{beamer}
\usetheme{Madrid}
\usecolortheme{whale}

\usepackage[utf8]{inputenc}
\usepackage{booktabs}
\usepackage{graphicx}

% Title page details
\title[RAG Semantic Poisoning]{Evaluating RAG Under Adversarial Data at Scale}
\subtitle{The Inverse Scaling Law of Semantic Poisoning \\ and the Dynamics of LLM-as-a-Judge}
\author[Nguyen et al.]{Nguyen Quoc Huy \and Bui Le Tan Dat \and Nguyen Dang Khoa \and Le Dinh Quy}
\institute[FPT University]{
  FPT University \\
  \medskip
  \textit{Supervisor: Le Thi Quynh Chi}
}
\date{\today}

\begin{document}

% Slide 1: Title
\begin{frame}
  \titlepage
\end{frame}

% Slide 2: Introduction
\begin{frame}{Introduction}
  \begin{itemize}
    \item \textbf{Background:} Retrieval-Augmented Generation (RAG) significantly reduces hallucinations by grounding Large Language Models (LLMs) in retrieved documents.
    \vspace{0.3cm}
    \item \textbf{The Vulnerability:} What happens when the retrieved documents are corrupted or semantically mutated? This creates a conflict between the LLM's parametric memory and the corrupted context.
    \vspace{0.3cm}
    \item \textbf{Objective:} Empirically evaluate RAG systems using 416 AI-generated semantic mutations to understand defensive behaviors and scaling laws.
  \end{itemize}
\end{frame}

% Slide 3: Methodology
\begin{frame}{Methodology}
  \begin{itemize}
    \item \textbf{Dataset:} 416 semantic mutations altering factual claims, numbers, and code.
    \begin{itemize}
      \item \textit{Public Data:} FastAPI documentation.
      \item \textit{Private Data:} Internal system documentation (no pretraining exposure).
    \end{itemize}
    \vspace{0.2cm}
    \item \textbf{Models Evaluated:} Llama 3.1 8B vs. 17B.
    \vspace{0.2cm}
    \item \textbf{Configurations:} Tested with and without defensive System Prompts.
    \vspace{0.2cm}
    \item \textbf{Evaluation Pipeline:} LLM-as-a-Judge (Llama 4 Scout 17B), cross-validated with human reviewers ($k > 0.73$).
  \end{itemize}
\end{frame}

% Slide 4: Evaluation Metrics (Labels)
\begin{frame}{Evaluation Metrics (Labels)}
  Every response was classified into one of four mutually exclusive labels:
  \vspace{0.4cm}
  \begin{enumerate}
    \item \textbf{Abstain (Safe):} The model recognizes contradictions and refuses to answer.
    \vspace{0.2cm}
    \item \textbf{Faithful (Fooled):} The model blindly repeats the poisoned context.
    \vspace{0.2cm}
    \item \textbf{Inconsistent (Violation):} The model silently corrects the error using its internal knowledge (violating RAG's groundedness principle).
    \vspace{0.2cm}
    \item \textbf{Hallucination:} The model fabricates completely unrelated information.
  \end{enumerate}
\end{frame}

% Slide 5: Key Finding 1 - The Inverse Scaling Law
\begin{frame}{Finding 1: The Inverse Scaling Law}
  \textbf{Hypothesis:} Larger models are safer and more resistant to attacks.\\
  \textbf{Reality:} Without a System Prompt, larger models are \textit{more vulnerable}.
  \vspace{0.3cm}
  \begin{table}
    \centering
    \begin{tabular}{lc}
      \toprule
      \textbf{Model (No Prompt)} & \textbf{Abstain Rate (Safety)} \\
      \midrule
      Llama 3.1 17B & 3.61\% \\
      Llama 3.1 8B  & 42.55\% \\
      \bottomrule
    \end{tabular}
  \end{table}
  \vspace{0.3cm}
  \textbf{Why?} The 17B model's superior reasoning allows it to seamlessly synthesize flawed logic (Faithful: 51.4\%) or silently fix it (Inconsistent: 44.7\%). The 8B model gets confused and safely abstains.
\end{frame}

% Slide 6: Key Finding 2 - The Parametric Shield
\begin{frame}{Finding 2: The Parametric Shield}
  What happens when the model evaluates \textbf{Private Data} (no pretraining exposure)?
  \vspace{0.3cm}
  \begin{itemize}
    \item The model loses its ``Parametric Shield'' (internal knowledge).
    \item \textbf{Inconsistent (self-correction)} rate plummeted to just 2.88\%.
    \item \textbf{Faithful (fooled)} rate jumped to 45.19\%.
  \end{itemize}
  \vspace{0.5cm}
  \begin{alertblock}{Conclusion}
    Enterprise RAG systems deployed on internal/private data are significantly more vulnerable to semantic poisoning than those using public data.
  \end{alertblock}
\end{frame}

% Slide 7: Key Finding 3 - Prompt Engineering Limits
\begin{frame}{Finding 3: Prompt Engineering is a Shield, Not a Cure}
  \begin{itemize}
    \item Adding a defensive System Prompt effectively acts as an immune response:
    \begin{itemize}
      \item 17B Abstain rate: 3.61\% $\rightarrow$ \textbf{60.10\%}
      \item 8B Abstain rate: 42.55\% $\rightarrow$ \textbf{63.70\%}
    \end{itemize}
    \vspace{0.3cm}
    \item However, \textbf{no configuration} achieved the $\ge 90\%$ Abstain safety threshold.
    \vspace{0.3cm}
    \item \textbf{Conclusion:} Lexical instructions alone cannot override deep semantic manipulation.
  \end{itemize}
\end{frame}

% Slide 8: Key Finding 4 - The Hallucination Myth
\begin{frame}{Finding 4: The Hallucination Myth}
  \begin{itemize}
    \item \textbf{Hallucinations} are often considered the biggest threat in LLMs.
    \vspace{0.2cm}
    \item Across over 2,000 automated evaluations, the hallucination rate \textbf{never exceeded 0.24\%} (effectively zero).
    \vspace{0.2cm}
    \item When poisoned, RAG models do not invent random facts. They exhibit a binary failure mode:
    \begin{enumerate}
      \item Passively trust the poison (Faithful).
      \item Rebel and silently fix the poison (Inconsistent).
    \end{enumerate}
    \vspace{0.2cm}
    \item \textbf{Shift in Focus:} Security research must pivot from measuring ``hallucinations'' to measuring ``faithfulness to poison.''
  \end{itemize}
\end{frame}

% Slide 9: Conclusion
\begin{frame}{Conclusion}
  \begin{itemize}
    \item \textbf{Scale vs. Safety:} RAG vulnerability scales \textit{inversely} with model size when left unprotected.
    \vspace{0.3cm}
    \item \textbf{Data Privacy Risk:} Private, internal data exposes RAG pipelines to much higher risks due to the absence of the Parametric Shield.
    \vspace{0.3cm}
    \item \textbf{Future Directions:} While prompt engineering helps, robust RAG systems will require structural and architectural defenses (e.g., GraphRAG, Controlled Noise Injection) to withstand semantic attacks.
  \end{itemize}
  \vspace{0.5cm}
  \begin{center}
    \Large \textbf{Thank You!}
  \end{center}
\end{frame}

\end{document}
